% Chapter 1: Introduction.

\section{Introduction}
\label{sec:intro}

\subsection{Problem Description}

A search engine answers a simple question---which documents contain this word
or phrase?---but the answer is only useful if it is cheap to obtain and
trustworthy when checked. This project builds such an engine for \emph{The
Complete Works of William Shakespeare}. The course asks for four pieces of
work: (1) count words and identify stop words; (2) build an inverted index with
stemming that excludes the stop words from part (1); (3) implement a query
program that accepts a word or phrase and returns document IDs; and (4) run
experiments that show how thresholds on a query affect the result. The bonus
question asks whether the design still works at the scale of 500,000 files and
400,000,000 distinct words.

The main difficulty is not any single algorithm but the coupling between the
stages. Stop-word removal changes position numbering; stemming changes the term
vocabulary; the query program must replay exactly the normalization used at
build time. A stop word that is left out of the index cannot be found later,
and a phrase that contains a stop word changes its meaning when that word is
dropped. The report therefore treats the pipeline as one specification rather
than four independent exercises.

\subsection{Purpose and Scope}

The purpose of this report is to document what was built, how it was tested,
and where the design stops working. Chapter~\ref{sec:design} specifies the
data structures and algorithms, Chapter~\ref{sec:results} reports the test
cases and measured results, and Chapter~\ref{sec:analysis} analyzes complexity,
compares the chosen structures with alternatives, and lists the limitations.
The report is deliberately evidence-oriented: every measured number comes from
a rerun on 2026-10-20, and the commands that reproduce the runs are collected
in Appendix~\ref{app:reproduce}. No measurement in this report was invented for
the write-up, and no number was recomputed from a different corpus or
configuration.

The scope is the single-machine, batch-built engine described here. Distributed
indexing, incremental updates, and Unicode tokenization are discussed as
future work but are not claimed as implemented features.

\subsection{Background}
\label{sec:background}

An \emph{inverted index} maps each term to a posting list that records the
documents containing the term and, in this project, the positions at which the
term occurs \cite{manning2008}. A word query intersects the posting lists of
its terms; a phrase query additionally checks that the positions are
consecutive. The two basic statistics are the term frequency $\mathrm{tf}(t)$,
the number of occurrences of term $t$ in the corpus, and the document frequency
$\mathrm{df}(t)$, the number of documents containing $t$. Their ratio separates
a word that is frequent because it occurs many times in one document from a
word that is spread across the collection.

\emph{Stop words} are terms with little discriminative value, such as function
words. There is no universal list: a word that is noise in one corpus can be
content in another. This project derives the list from the corpus itself and
checks the derivation on a second corpus. \emph{Stemming} reduces inflected
forms to a common stem so that, for example, \code{agree}, \code{agreed}, and
\code{agreement} share an index entry \cite{porter1980}. Stemming improves
recall but can over-conflate unrelated words, so the implementation is checked
against an independent implementation on the full vocabulary.

The dictionary that maps terms to posting lists is itself a data structure
problem. A balanced search tree, a hash table, a splay tree, and a disk-based
B+ tree offer different trade-offs between lookup time, ordered traversal, and
memory \cite{cormen2022}. The query thresholds are a separate design axis:
\code{max-df} and \code{min-df} filter query terms, \code{min-match} changes the
Boolean condition, and \code{top-x} truncates the ranked output. The experiment
in Chapter~\ref{sec:results} shows that these controls are not interchangeable.

\subsection{Corpus Collection}

\subsubsection{Corpus Source}
\label{sec:corpus}

The main corpus is \emph{The Complete Works of William Shakespeare} from the
MIT Shakespeare site \cite{shakespeare} and its official GitHub mirror. The
mirror contains three overlapping copies of the works: the current lowercase
play directories, an older uppercase classification directory, and Moby plain
text files. To avoid mixing editions, the site's own index page is used as the
authoritative table of contents. This yields 195 documents: 37 plays, 154
sonnets, and 4 long poems, totaling 152,632 lines and 930,911 words in
4,919,371 bytes of source text.

A second corpus checks that the engine is not tied to Shakespeare. It contains
three IETF RFCs---1035 (DNS), 5321 (SMTP), and 7231 (HTTP semantics)---stored
as plain text: 3 documents, 9,946 lines, and 81,848 whitespace-separated words.
The same C++ binary processes both corpora; only the manifest and the text
layer change.

\subsubsection{Preprocessing Overview}
\label{sec:preprocess}

\code{tools/extract.py} converts each HTML page into a text layer with one
semantic unit per line, and writes \code{manifest.tsv} as the document catalog
and path map. Line markers identify headings such as \code{[ACT I]}, stage
directions such as \code{(Exit Ghost)}, and speaker labels such as
\code{@BERNARDO}; the tokenizer ignores these lines by default. The marker
must match at both the first and the last character. A line that only starts
with \code{[} but continues as dialogue is body text, not a heading. This rule
is essential for the corpus counts: an earlier comparison script tested only
the first character, treated lines such as \code{[To HELENA] The best wishes
...} as headings, and undercounted the token stream by 4,340 tokens.

The detailed token rules are part of the specification in
Section~\ref{sec:tokenization}; the corpus-level view is enough here. All
tokens are lowercased. The default configuration counts body lines only and
produces 866,626 tokens. The flags \code{--keep-headings},
\code{--keep-speakers}, and \code{--keep-stage} can retain each marker class
when a different counting convention is needed.

\subsubsection{Special Cases}
\label{sec:preprocess-analysis}

Three corpus defects required explicit handling. First, the Moby edition writes
610 bracketed stage directions; 57 are whole-line forms such as
\code{[Sings]}, and 553 are line-start fragments followed by dialogue, such as
\code{[To HELENA] The best wishes ...}. The extractor relocates these to
parenthesized stage lines. Eight mid-line bracketed forms are editorial
insertions or missing-text markers and are left in place. The net effect is a
reduction of 849 tokens, which matches the fragment counts exactly. Second,
\code{Poetry/elegy.html} contains the malformed tag \code{</TITLE} without a
closing angle bracket; the HTML parser never leaves the head and extracts no
text. Treating the start of \code{<body>} as the end of the head fixes the
file. Third, \code{2henryiv.html} contains a source artifact,
\code{<h3>None</h3>}, which is discarded rather than turned into a word.

The text layer and the token stream are checked against an independent Python
implementation, so a preprocessing error is visible as a per-document
difference rather than as a plausible but wrong total. The individual test
cases are listed in Section~\ref{sec:test-cases}.

\subsection{System Overview}
\label{sec:architecture}

Figure~\ref{fig:pipeline} shows the data flow. On the offline side, the HTML
corpus becomes a text layer, the tokenizer produces a token stream, and the
stream feeds two branches. The counting branch writes \code{wordcount.csv} and
derives the stop-word list. The indexing branch applies the same stop-word
filter and the optional Porter stemmer, inserts each term into the dictionary,
and serializes the postings to \code{index.bin}; an optional B+ tree sidecar
stores the term-to-posting mapping for lazy lookup. On the online side, the
query text goes through the same tokenizer, stop-word filter, and stemmer, then
looks up the dictionary, intersects posting lists, verifies positions for
phrases, applies the thresholds, and returns document IDs and titles.

% figures/pipeline.tex -- system data flow (offline build and online query).
% Pure TikZ diagram: stage names, file names, and mechanisms only; no measured
% values. The main document writes % figures/pipeline.tex -- system data flow (offline build and online query).
% Pure TikZ diagram: stage names, file names, and mechanisms only; no measured
% values. The main document writes % figures/pipeline.tex -- system data flow (offline build and online query).
% Pure TikZ diagram: stage names, file names, and mechanisms only; no measured
% values. The main document writes \input{figures/pipeline}.
\begin{figure}[!htbp]
\centering
\resizebox{\linewidth}{!}{%
\begin{tikzpicture}[
  font=\footnotesize,
  >=Stealth,
  x=1cm, y=1cm,
  data/.style={draw=black!55, rounded corners=2pt, line width=0.5pt,
               fill=black!4, inner sep=3pt, align=center, text width=2.25cm},
  proc/.style={draw=mseblue, rounded corners=2pt, line width=0.6pt,
               fill=mseblue!6, text=mseblue, inner sep=3pt, align=center,
               text width=2.25cm},
  note/.style={font=\scriptsize, text=msegray, align=center},
  group/.style={draw=black!18, rounded corners=6pt, line width=0.5pt,
                dashed, fill=black!1},
  grouplabel/.style={font=\scriptsize, text=msegray, fill=white,
                     inner sep=1.5pt},
  arr/.style={->, line width=0.6pt, draw=black!65},
  darr/.style={->, dashed, line width=0.7pt, draw=mseblue},
  fork/.style={circle, fill=black!60, inner sep=0pt, minimum size=2.4pt},
]

% ---------------------------------------------------------------- groups
\draw[group] (-0.15,2.65) rectangle (15.55,7.85);
\node[grouplabel] at (1.65,7.85) {Offline build};
\draw[group] (-0.15,0.18) rectangle (15.55,1.72);
\node[grouplabel] at (1.65,1.72) {Online query};

% ---------------------------------------------------------------- corpus
\node[data] (raw) at (1.75,7.15) {\texttt{data/raw/*.html}};
\node[proc] (extract) at (4.55,7.15) {\texttt{extract.py}};
\node[data] (text) at (7.85,7.15) {\texttt{data/text/*.txt}\\[1pt]\texttt{manifest.tsv}};
\node[proc] (tokenizer) at (11.10,7.15) {\texttt{tokenizer}\\[1pt]{\scriptsize line by line}};
\node[data] (tokens) at (13.85,7.15) {token stream\\[1pt]{\scriptsize lowercase, letters}};

\draw[arr] (raw) -- (extract);
\draw[arr] (extract) -- (text);
\draw[arr] (text) -- (tokenizer);
\draw[arr] (tokenizer) -- (tokens);

\node[note] at (7.55,6.38)
  {markers: \texttt{@SPEAKER} / \texttt{(stage)} / \texttt{[ACT]}};

% ---------------------------------------------------------------- fork: count
\coordinate (fork) at (13.85,6.05);
\draw[arr] (tokens.south) -- (fork);
\node[fork] at (fork) {};
\draw[arr] (fork) -- (fork -| 0.85,0) -- (0.85,3.62);

\draw[arr] (fork) -- (13.85,5.62);
\node[proc] (count) at (13.85,5.15) {\texttt{count}\\[1pt]{\scriptsize tf, df}};
\node[data] (wc) at (11.15,5.15) {\texttt{wordcount.csv}};
\node[proc] (rule) at (8.20,5.15) {three-feature\\rule\\[1pt]{\scriptsize df, tf, entropy}};
\node[data] (sw) at (5.00,5.15) {\texttt{data/}\\[-2pt]\texttt{stopwords.txt}};
\draw[arr] (count) -- (wc);
\draw[arr] (wc) -- (rule);
\draw[arr] (rule) -- (sw);

\draw[darr] (sw.south) -- (5.00,4.28) -- (1.90,4.28) -- (1.90,3.62);
\node[note] at (3.45,4.28) {stop-word list};

% ---------------------------------------------------------------- index build
\node[proc] (stop) at (1.45,3.25) {\texttt{stop filter}\\[1pt]{\scriptsize 39 words}};
\node[proc] (stem) at (4.05,3.25) {\texttt{Porter}\\[1pt]{\scriptsize stemmer}};
\node[proc] (dict) at (6.65,3.25) {\texttt{dictionary}\\[1pt]{\scriptsize splay/hash/map}};
\node[proc] (post) at (9.25,3.25) {\texttt{postings}\\[1pt]{\scriptsize docID, tf, pos}};
\node[proc] (ser) at (11.85,3.25) {\texttt{serialize}\\[1pt]{\scriptsize v1 / v2}};
\node[data] (idx) at (14.35,3.25) {\texttt{index.bin}\\[1pt]\texttt{index.bpt}\\[1pt]{\scriptsize sidecar}};

\draw[arr] (stop) -- (stem);
\draw[arr] (stem) -- (dict);
\draw[arr] (dict) -- (post);
\draw[arr] (post) -- (ser);
\draw[arr] (ser) -- (idx);

% ---------------------------------------------------------------- query side
\node[proc] (query) at (1.45,0.95) {\texttt{mse query}\\[1pt]{\scriptsize word/phrase}};
\node[proc] (norm) at (4.05,0.95) {same pipeline\\[1pt]{\scriptsize tokenize, stop, stem}};
\node[proc] (lookup) at (6.65,0.95) {dictionary\\lookup\\[1pt]{\scriptsize $O(\log T)$}};
\node[proc] (intersect) at (9.25,0.95) {posting\\intersection\\[1pt]{\scriptsize position check}};
\node[proc] (thresh) at (11.85,0.95) {thresholds\\and ranking\\[1pt]{\scriptsize max-df/min-df}\\[1pt]{\scriptsize min-match/top-x}};
\node[data] (hits) at (14.35,0.95) {document IDs\\[1pt]{\scriptsize ranked list}};

\draw[arr] (query) -- (norm);
\draw[arr] (norm) -- (lookup);
\draw[arr] (lookup) -- (intersect);
\draw[arr] (intersect) -- (thresh);
\draw[arr] (thresh) -- (hits);

\draw[darr] (idx.south) -- (14.35,2.10) -- (6.65,2.10) -- (lookup.north);
\node[note, fill=white, inner sep=2pt] at (10.45,2.10)
  {\texttt{loadIndex}: full term table\\[1pt]
   \texttt{index.bpt} + \texttt{LazyIndexReader}: $O(h)$ page reads};

\end{tikzpicture}%
}
\caption{System data flow. The offline side extracts the HTML corpus into a
text layer, tokenizes it, removes stop words, applies Porter stemming, and
writes the inverted index; word-frequency counting and stop-word derivation
form a bypass. The online side reuses the same normalization pipeline, then
performs dictionary lookup, posting intersection and position verification,
applies the four thresholds, and returns document IDs. The dashed path is the
B+ tree sidecar written by \code{--dict-file}; a query reads $O(h)$ pages,
where $h$ is the tree height.}
\label{fig:pipeline}
\end{figure}
.
\begin{figure}[!htbp]
\centering
\resizebox{\linewidth}{!}{%
\begin{tikzpicture}[
  font=\footnotesize,
  >=Stealth,
  x=1cm, y=1cm,
  data/.style={draw=black!55, rounded corners=2pt, line width=0.5pt,
               fill=black!4, inner sep=3pt, align=center, text width=2.25cm},
  proc/.style={draw=mseblue, rounded corners=2pt, line width=0.6pt,
               fill=mseblue!6, text=mseblue, inner sep=3pt, align=center,
               text width=2.25cm},
  note/.style={font=\scriptsize, text=msegray, align=center},
  group/.style={draw=black!18, rounded corners=6pt, line width=0.5pt,
                dashed, fill=black!1},
  grouplabel/.style={font=\scriptsize, text=msegray, fill=white,
                     inner sep=1.5pt},
  arr/.style={->, line width=0.6pt, draw=black!65},
  darr/.style={->, dashed, line width=0.7pt, draw=mseblue},
  fork/.style={circle, fill=black!60, inner sep=0pt, minimum size=2.4pt},
]

% ---------------------------------------------------------------- groups
\draw[group] (-0.15,2.65) rectangle (15.55,7.85);
\node[grouplabel] at (1.65,7.85) {Offline build};
\draw[group] (-0.15,0.18) rectangle (15.55,1.72);
\node[grouplabel] at (1.65,1.72) {Online query};

% ---------------------------------------------------------------- corpus
\node[data] (raw) at (1.75,7.15) {\texttt{data/raw/*.html}};
\node[proc] (extract) at (4.55,7.15) {\texttt{extract.py}};
\node[data] (text) at (7.85,7.15) {\texttt{data/text/*.txt}\\[1pt]\texttt{manifest.tsv}};
\node[proc] (tokenizer) at (11.10,7.15) {\texttt{tokenizer}\\[1pt]{\scriptsize line by line}};
\node[data] (tokens) at (13.85,7.15) {token stream\\[1pt]{\scriptsize lowercase, letters}};

\draw[arr] (raw) -- (extract);
\draw[arr] (extract) -- (text);
\draw[arr] (text) -- (tokenizer);
\draw[arr] (tokenizer) -- (tokens);

\node[note] at (7.55,6.38)
  {markers: \texttt{@SPEAKER} / \texttt{(stage)} / \texttt{[ACT]}};

% ---------------------------------------------------------------- fork: count
\coordinate (fork) at (13.85,6.05);
\draw[arr] (tokens.south) -- (fork);
\node[fork] at (fork) {};
\draw[arr] (fork) -- (fork -| 0.85,0) -- (0.85,3.62);

\draw[arr] (fork) -- (13.85,5.62);
\node[proc] (count) at (13.85,5.15) {\texttt{count}\\[1pt]{\scriptsize tf, df}};
\node[data] (wc) at (11.15,5.15) {\texttt{wordcount.csv}};
\node[proc] (rule) at (8.20,5.15) {three-feature\\rule\\[1pt]{\scriptsize df, tf, entropy}};
\node[data] (sw) at (5.00,5.15) {\texttt{data/}\\[-2pt]\texttt{stopwords.txt}};
\draw[arr] (count) -- (wc);
\draw[arr] (wc) -- (rule);
\draw[arr] (rule) -- (sw);

\draw[darr] (sw.south) -- (5.00,4.28) -- (1.90,4.28) -- (1.90,3.62);
\node[note] at (3.45,4.28) {stop-word list};

% ---------------------------------------------------------------- index build
\node[proc] (stop) at (1.45,3.25) {\texttt{stop filter}\\[1pt]{\scriptsize 39 words}};
\node[proc] (stem) at (4.05,3.25) {\texttt{Porter}\\[1pt]{\scriptsize stemmer}};
\node[proc] (dict) at (6.65,3.25) {\texttt{dictionary}\\[1pt]{\scriptsize splay/hash/map}};
\node[proc] (post) at (9.25,3.25) {\texttt{postings}\\[1pt]{\scriptsize docID, tf, pos}};
\node[proc] (ser) at (11.85,3.25) {\texttt{serialize}\\[1pt]{\scriptsize v1 / v2}};
\node[data] (idx) at (14.35,3.25) {\texttt{index.bin}\\[1pt]\texttt{index.bpt}\\[1pt]{\scriptsize sidecar}};

\draw[arr] (stop) -- (stem);
\draw[arr] (stem) -- (dict);
\draw[arr] (dict) -- (post);
\draw[arr] (post) -- (ser);
\draw[arr] (ser) -- (idx);

% ---------------------------------------------------------------- query side
\node[proc] (query) at (1.45,0.95) {\texttt{mse query}\\[1pt]{\scriptsize word/phrase}};
\node[proc] (norm) at (4.05,0.95) {same pipeline\\[1pt]{\scriptsize tokenize, stop, stem}};
\node[proc] (lookup) at (6.65,0.95) {dictionary\\lookup\\[1pt]{\scriptsize $O(\log T)$}};
\node[proc] (intersect) at (9.25,0.95) {posting\\intersection\\[1pt]{\scriptsize position check}};
\node[proc] (thresh) at (11.85,0.95) {thresholds\\and ranking\\[1pt]{\scriptsize max-df/min-df}\\[1pt]{\scriptsize min-match/top-x}};
\node[data] (hits) at (14.35,0.95) {document IDs\\[1pt]{\scriptsize ranked list}};

\draw[arr] (query) -- (norm);
\draw[arr] (norm) -- (lookup);
\draw[arr] (lookup) -- (intersect);
\draw[arr] (intersect) -- (thresh);
\draw[arr] (thresh) -- (hits);

\draw[darr] (idx.south) -- (14.35,2.10) -- (6.65,2.10) -- (lookup.north);
\node[note, fill=white, inner sep=2pt] at (10.45,2.10)
  {\texttt{loadIndex}: full term table\\[1pt]
   \texttt{index.bpt} + \texttt{LazyIndexReader}: $O(h)$ page reads};

\end{tikzpicture}%
}
\caption{System data flow. The offline side extracts the HTML corpus into a
text layer, tokenizes it, removes stop words, applies Porter stemming, and
writes the inverted index; word-frequency counting and stop-word derivation
form a bypass. The online side reuses the same normalization pipeline, then
performs dictionary lookup, posting intersection and position verification,
applies the four thresholds, and returns document IDs. The dashed path is the
B+ tree sidecar written by \code{--dict-file}; a query reads $O(h)$ pages,
where $h$ is the tree height.}
\label{fig:pipeline}
\end{figure}
.
\begin{figure}[!htbp]
\centering
\resizebox{\linewidth}{!}{%
\begin{tikzpicture}[
  font=\footnotesize,
  >=Stealth,
  x=1cm, y=1cm,
  data/.style={draw=black!55, rounded corners=2pt, line width=0.5pt,
               fill=black!4, inner sep=3pt, align=center, text width=2.25cm},
  proc/.style={draw=mseblue, rounded corners=2pt, line width=0.6pt,
               fill=mseblue!6, text=mseblue, inner sep=3pt, align=center,
               text width=2.25cm},
  note/.style={font=\scriptsize, text=msegray, align=center},
  group/.style={draw=black!18, rounded corners=6pt, line width=0.5pt,
                dashed, fill=black!1},
  grouplabel/.style={font=\scriptsize, text=msegray, fill=white,
                     inner sep=1.5pt},
  arr/.style={->, line width=0.6pt, draw=black!65},
  darr/.style={->, dashed, line width=0.7pt, draw=mseblue},
  fork/.style={circle, fill=black!60, inner sep=0pt, minimum size=2.4pt},
]

% ---------------------------------------------------------------- groups
\draw[group] (-0.15,2.65) rectangle (15.55,7.85);
\node[grouplabel] at (1.65,7.85) {Offline build};
\draw[group] (-0.15,0.18) rectangle (15.55,1.72);
\node[grouplabel] at (1.65,1.72) {Online query};

% ---------------------------------------------------------------- corpus
\node[data] (raw) at (1.75,7.15) {\texttt{data/raw/*.html}};
\node[proc] (extract) at (4.55,7.15) {\texttt{extract.py}};
\node[data] (text) at (7.85,7.15) {\texttt{data/text/*.txt}\\[1pt]\texttt{manifest.tsv}};
\node[proc] (tokenizer) at (11.10,7.15) {\texttt{tokenizer}\\[1pt]{\scriptsize line by line}};
\node[data] (tokens) at (13.85,7.15) {token stream\\[1pt]{\scriptsize lowercase, letters}};

\draw[arr] (raw) -- (extract);
\draw[arr] (extract) -- (text);
\draw[arr] (text) -- (tokenizer);
\draw[arr] (tokenizer) -- (tokens);

\node[note] at (7.55,6.38)
  {markers: \texttt{@SPEAKER} / \texttt{(stage)} / \texttt{[ACT]}};

% ---------------------------------------------------------------- fork: count
\coordinate (fork) at (13.85,6.05);
\draw[arr] (tokens.south) -- (fork);
\node[fork] at (fork) {};
\draw[arr] (fork) -- (fork -| 0.85,0) -- (0.85,3.62);

\draw[arr] (fork) -- (13.85,5.62);
\node[proc] (count) at (13.85,5.15) {\texttt{count}\\[1pt]{\scriptsize tf, df}};
\node[data] (wc) at (11.15,5.15) {\texttt{wordcount.csv}};
\node[proc] (rule) at (8.20,5.15) {three-feature\\rule\\[1pt]{\scriptsize df, tf, entropy}};
\node[data] (sw) at (5.00,5.15) {\texttt{data/}\\[-2pt]\texttt{stopwords.txt}};
\draw[arr] (count) -- (wc);
\draw[arr] (wc) -- (rule);
\draw[arr] (rule) -- (sw);

\draw[darr] (sw.south) -- (5.00,4.28) -- (1.90,4.28) -- (1.90,3.62);
\node[note] at (3.45,4.28) {stop-word list};

% ---------------------------------------------------------------- index build
\node[proc] (stop) at (1.45,3.25) {\texttt{stop filter}\\[1pt]{\scriptsize 39 words}};
\node[proc] (stem) at (4.05,3.25) {\texttt{Porter}\\[1pt]{\scriptsize stemmer}};
\node[proc] (dict) at (6.65,3.25) {\texttt{dictionary}\\[1pt]{\scriptsize splay/hash/map}};
\node[proc] (post) at (9.25,3.25) {\texttt{postings}\\[1pt]{\scriptsize docID, tf, pos}};
\node[proc] (ser) at (11.85,3.25) {\texttt{serialize}\\[1pt]{\scriptsize v1 / v2}};
\node[data] (idx) at (14.35,3.25) {\texttt{index.bin}\\[1pt]\texttt{index.bpt}\\[1pt]{\scriptsize sidecar}};

\draw[arr] (stop) -- (stem);
\draw[arr] (stem) -- (dict);
\draw[arr] (dict) -- (post);
\draw[arr] (post) -- (ser);
\draw[arr] (ser) -- (idx);

% ---------------------------------------------------------------- query side
\node[proc] (query) at (1.45,0.95) {\texttt{mse query}\\[1pt]{\scriptsize word/phrase}};
\node[proc] (norm) at (4.05,0.95) {same pipeline\\[1pt]{\scriptsize tokenize, stop, stem}};
\node[proc] (lookup) at (6.65,0.95) {dictionary\\lookup\\[1pt]{\scriptsize $O(\log T)$}};
\node[proc] (intersect) at (9.25,0.95) {posting\\intersection\\[1pt]{\scriptsize position check}};
\node[proc] (thresh) at (11.85,0.95) {thresholds\\and ranking\\[1pt]{\scriptsize max-df/min-df}\\[1pt]{\scriptsize min-match/top-x}};
\node[data] (hits) at (14.35,0.95) {document IDs\\[1pt]{\scriptsize ranked list}};

\draw[arr] (query) -- (norm);
\draw[arr] (norm) -- (lookup);
\draw[arr] (lookup) -- (intersect);
\draw[arr] (intersect) -- (thresh);
\draw[arr] (thresh) -- (hits);

\draw[darr] (idx.south) -- (14.35,2.10) -- (6.65,2.10) -- (lookup.north);
\node[note, fill=white, inner sep=2pt] at (10.45,2.10)
  {\texttt{loadIndex}: full term table\\[1pt]
   \texttt{index.bpt} + \texttt{LazyIndexReader}: $O(h)$ page reads};

\end{tikzpicture}%
}
\caption{System data flow. The offline side extracts the HTML corpus into a
text layer, tokenizes it, removes stop words, applies Porter stemming, and
writes the inverted index; word-frequency counting and stop-word derivation
form a bypass. The online side reuses the same normalization pipeline, then
performs dictionary lookup, posting intersection and position verification,
applies the four thresholds, and returns document IDs. The dashed path is the
B+ tree sidecar written by \code{--dict-file}; a query reads $O(h)$ pages,
where $h$ is the tree height.}
\label{fig:pipeline}
\end{figure}


Keeping the two branches on one tokenizer is a deliberate constraint. It makes
the stop-word list, the stemmer, and the query program share a single
definition of a term, and it lets the index header record the processing flags
so that a query cannot silently use a different normalization.

\subsection{Report Organization}

Chapter~\ref{sec:design} gives the data-structure and algorithm specification,
including pseudo-code for the main operations. Chapter~\ref{sec:results}
presents the test environment, the required test-case table, and the measured
results for the four tasks and the bonus experiments.
Chapter~\ref{sec:analysis} analyzes time and space complexity, compares the
chosen structures with alternatives, generalizes to the RFC corpus, and lists
limitations and improvements. Appendix~\ref{app:source} records the source-code
organization and comment coverage, Appendix~\ref{app:reproduce} gives the
reproduction commands, and Appendix~\ref{app:format} documents the data
formats. The Author List, Declaration, and Signatures follow the appendices.
